\documentclass[10pt,a4paper]{amsart}
\usepackage[margin=19mm]{geometry}
\usepackage{amsmath,amssymb,mathtools}
\usepackage[scr=rsfs,cal=euler]{mathalfa}
\usepackage{xfrac}
\usepackage[unicode,hidelinks,bookmarks=false]{hyperref}
\newcommand{\ie}{i.e.\ }
\newcommand{\cf}{cf.\ }
\newcommand{\resp}{respectively\ }
\newcommand{\Subsection}{\subsection}
\providecommand{\nobreakdash}{\nobreak\hbox{-}}
\setlength{\emergencystretch}{2em}
\multlinegap0pt
\usepackage{amssymb}
\usepackage{mathtools}
\usepackage{algorithm}\usepackage{algpseudocode}
\usepackage{tikz}
\usepackage[shortlabels]{enumitem}
\newtheorem{lemma}{Lemma}
\usepackage{cleveref}
\crefname{lemma}{Lemma}{Lemmas}
\newcommand{\cD}{\mathcal D}
\newcommand{\norm}[1]{\lVert #1\rVert}
\title{A deterministic $(1+\varepsilon)^n$ approximation for the permanent of a nonnegative matrix}
\author{Dingding Dong, Vishesh Jain}
\date{Source: arXiv:2609.11049v1 (CC BY 4.0)}
\begin{document}
\maketitle

\noindent\emph{Source and licence.} A deterministic $(1+\varepsilon)^n$ approximation for the permanent of a nonnegative matrix. Version arXiv:2609.11049v1 (CC BY 4.0), distributed by arXiv under the Creative Commons Attribution 4.0 International licence, http://creativecommons.org/licenses/by/4.0/, as recorded in the arXiv licence metadata for this exact version and shown on the abstract page https://arxiv.org/abs/2609.11049v1. Full licence notice: \texttt{licenses/arxiv-2609.11049-CC-BY-4.0.txt}.\par\smallskip

\noindent\textit{Editorial scope.} Counted unit: the article's lemma lem:analytic-derivatives (source0.tex lines 1558--1574). The article's complete original proof of it is delivered with it (source0.tex lines 1576--1726, one \texttt{proof} environment of 151 lines). No mathematics is written, repaired or re-derived: the statement and the proof are the article's own lines, in the article's own notation, and no retained line is substituted or re-flowed.  Retained uncounted as the source-local context the proof needs: lines 1344--1354; lines 1347--1349; lines 1491--1498.  Nothing was omitted from any retained span, so the package carries no omission record.  No retained line contains a citation, so the wrapper carries no bibliography and no bibliography entry.  No retained line contains a figure environment, an image-inclusion command or drawing code, so no image is delivered and the package declares no asset dependency.  Editorial formatting only, in addition: an AMS \texttt{amsart} preamble that restates the article's own declarations and macros used by the retained lines, namely the environments \texttt{lemma} on separate counters, together with the standard packages those lines need.  The retained environments print in this document as Lemma 1 in order of appearance, whereas the article prints the same items with the numbering of its own full document.  The complete native source of the article as distributed in the arXiv e-print for this exact version is bundled unchanged as \texttt{sources/209982/source0.tex}.
\begin{lemma}\label{lem:path-tree-resolvent}
For every induced subgraph $H$ of $G$ and $v\in V(H)$, let $B$ and
$r$ be as above. Then
\begin{equation}\label{eq:path-tree-resolvent}
 R_{H,v}(sw)=\bigl[(I+sB^2)^{-1}\bigr]_{rr}
\end{equation}
whenever both sides are defined. If $w=e^y$ with $y\in\cD_K$, then
\begin{equation*}
 \norm{B}_2\le2\sqrt K.
\end{equation*}
\end{lemma}

\begin{equation}
 R_{H,v}(sw)=\bigl[(I+sB^2)^{-1}\bigr]_{rr}
\end{equation}

\begin{equation}\label{eq:interpolation-contour-bounds}
 \begin{aligned}
  \frac{z_\star}{R}&=\frac{8K}{8K+1}=q,\\
  |\tau(z)|&\le\frac{\rho R}{1-R}
             =2+\frac1{4K}\le\frac94
             \qquad(|z|\le R).
 \end{aligned}
\end{equation}

\begin{lemma}\label{lem:analytic-derivatives}
For every induced subgraph $H$ of $G$,
$Z_H(\tau(z)w)\ne0$ for $|z|<1$. In particular, there is a
holomorphic logarithm
\[
 F_y(z)=\log Z_G(\tau(z)e^y)
 \qquad(|z|<1)
\]
with $F_y(0)=0$. For every edge $e$, every
$\sigma\in\{-1,1\}^E$, and $|z|\le R$,
\begin{align}
 |\partial_eF_y(z)|&\le9e^{y_e},
 \label{eq:analytic-gradient-bound}\\
 |D_\sigma\partial_eF_y(z)|&\le400K e^{y_e}.
 \label{eq:analytic-signed-hessian-bound}
\end{align}
\end{lemma}

\begin{proof}
For $|z|<1$,
\[
 \operatorname{Re}\frac{z}{1-z}
 =\frac12\left(\frac{1-|z|^2}{|1-z|^2}-1\right)>-\frac12,
\]
so $\operatorname{Re}\tau(z)>-1/(8K)$. Let $B$ be the matrix
associated with any path tree $\mathcal T(H,v)$.
By \cref{lem:path-tree-resolvent}, $\norm B_2\le2\sqrt K$,
so every eigenvalue $a$ of $B^2$ lies in $[0,4K]$. Hence
\[
 \operatorname{Re}(1+\tau(z)a)
 \ge1-\frac{a}{8K}\ge\frac12.
\]
Since $B^2$ is real symmetric,
\begin{equation*}
 \norm{(I+\tau(z)B^2)^{-1}}_2
 =\max_{a\in\operatorname{spec}(B^2)}
       \frac1{|1+\tau(z)a|}
 \le2
 \qquad(|z|<1).
\end{equation*}

We next prove that $Z_H(\tau(z)w)\ne0$ for every induced subgraph
$H$, by induction on $|V(H)|$. The empty graph has partition
function $1$. For nonempty $H$, choose $v\in V(H)$ and let $B$
be the matrix of $\mathcal T(H,v)$, with root $r$.
By \cref{eq:path-tree-resolvent},
\[
 Z_H(sw)\bigl[(I+sB^2)^{-1}\bigr]_{rr}=Z_{H-v}(sw)
\]
as an identity of rational functions in $s$. By the inverse bound above,
the inverse exists at $s=\tau(z)$, so the identity holds there.
Its right-hand side is nonzero by induction, and therefore
$Z_H(\tau(z)w)\ne0$.

The unit disk is simply connected, so $Z_G(\tau(z)w)$ admits a
holomorphic logarithm there. Since $Z_G(\tau(0)w)=1$, we may choose
it to satisfy $F_y(0)=0$. By continuity and compactness,
$Z_G(\tau(z)e^{y'})$ remains nonzero on a common disk containing
$|z|\le R$ for all $y'$ sufficiently close to $y$. On this disk,
choose each logarithm to satisfy $F_{y'}(0)=0$. These logarithms
depend smoothly on $y'$, so we may differentiate with respect
to $y$, using the usual rule
\[
 \partial_eF_y(z)
 =\frac{\partial_e Z_G(\tau(z)e^y)}{Z_G(\tau(z)e^y)}.
\]

For any induced subgraph $H$ and vertex $v\in V(H)$, the path-tree
representation \cref{eq:path-tree-resolvent} and the inverse norm bound
above give
\begin{equation*}
 |R_{H,v}(\tau(z)w)|
 =\left|\bigl[(I+\tau(z)B^2)^{-1}\bigr]_{rr}\right|
 \le2
 \qquad(|z|<1).
\end{equation*}
For $e=uv$, differentiating the partition function gives
\[
 \partial_eF_y(z)
 =\tau(z)w_e\frac{Z_{G-u-v}(\tau(z)w)}{Z_G(\tau(z)w)}
 =\tau(z)w_eR_{G,u}(\tau(z)w)R_{G-u,v}(\tau(z)w).
\]
Since $|\tau(z)|\le9/4$ on $|z|\le R$ by
\cref{eq:interpolation-contour-bounds}, it follows that
\[
 |\partial_eF_y(z)|
 \le\frac94w_e\cdot2\cdot2=9w_e
 \qquad(|z|\le R).
\]
This proves \cref{eq:analytic-gradient-bound}.

To bound $D_\sigma\partial_eF_y(z)$, we first differentiate the
matrix expression for $R_{H,v}$. Write $B_\sigma=D_\sigma B$,
where the derivative is taken entrywise. If the edge joining
$p$ and $p'$ corresponds to $f\in E(H)$, then
\[
 (B_\sigma)_{pp'}
 =D_\sigma(e^{y_f/2})
 =\frac{\sigma_f}{2}B_{pp'}.
\]
As in the path-tree proof, write $B=P+P^{\mathsf T}$, where
$P_{pp'}=B_{pp'}$ when $p$ is the parent of $p'$, and all other
entries of $P$ are zero. Let $P_\sigma=D_\sigma P$, so that
\[
 B_\sigma=P_\sigma+P_\sigma^{\mathsf T}.
\]
As before, $P_\sigma P_\sigma^{\mathsf T}$ is diagonal.
Its diagonal entries are one quarter of those of
$PP^{\mathsf T}$, because each entry of $P_\sigma$ has half
the absolute value of the corresponding entry of $P$. Hence
\[
 P_\sigma P_\sigma^{\mathsf T}
 =\frac14PP^{\mathsf T}\preceq\frac K4I.
\]
It follows that
\begin{equation*}
 \norm{B_\sigma}_2
 \le\norm{P_\sigma}_2+\norm{P_\sigma^{\mathsf T}}_2
 \le\sqrt K.
\end{equation*}

Differentiating the identity
$(I+\tau(z)B^2)(I+\tau(z)B^2)^{-1}=I$ gives
\[
 D_\sigma(I+\tau(z)B^2)^{-1}
 =-(I+\tau(z)B^2)^{-1}
   \tau(z)(B_\sigma B+BB_\sigma)
   (I+\tau(z)B^2)^{-1}.
\]
The ratio $R_{H,v}(\tau(z)w)$ is the $rr$ entry of this inverse.
Therefore, using the bounds on the inverse, $\tau(z)$, $B$,
and $B_\sigma$, we obtain
\begin{align*}
 |D_\sigma R_{H,v}(\tau(z)w)|
 &\le
 2|\tau(z)|\norm{(I+\tau(z)B^2)^{-1}}_2^2
   \norm B_2\norm{B_\sigma}_2\\
 &\le2\cdot\frac94\cdot2^2\cdot2\sqrt K\cdot\sqrt K
 =36K
 \qquad(|z|\le R).
\end{align*}

Finally, for $e=uv$, let
\[
 R_1=R_{G,u}(\tau(z)w),
 \qquad
 R_2=R_{G-u,v}(\tau(z)w).
\]
Since $\partial_eF_y(z)=\tau(z)w_eR_1R_2$ and
$D_\sigma w_e=\sigma_ew_e$, the product rule gives
\[
 D_\sigma\partial_eF_y(z)
 =\tau(z)w_e\bigl(
   \sigma_eR_1R_2+(D_\sigma R_1)R_2+R_1(D_\sigma R_2)
   \bigr).
\]
On $|z|\le R$, each ratio has absolute value at most $2$
by the bounds established above, and each of its signed derivatives
has absolute value at most $36K$ by
the derivative bound above. Thus,
\[
 |D_\sigma\partial_eF_y(z)|
 \le\frac94w_e\bigl(4+36K\cdot2+2\cdot36K\bigr)
 =(9+324K)w_e
 \le400K w_e,
\]
where the last inequality uses $K\ge1$. This proves
\cref{eq:analytic-signed-hessian-bound}.
\end{proof}

\end{document}
